\begin{afterword}
\summaryhead

The author recalls knowing, at nine, that light, sound and matter were waves, and finding out years later, after three days with the wave equation in the \textsc{Feynman Lectures on Physics}, that the younger self had had no idea what ``wave'' means to a physicist. School, the post argues, rewards saying the right words. A student who answers ``Waves!'' has not thereby made a true statement: the syllables are not a hypothesis, only verbal behavior. A student asked why the far side of a metal plate is warmer, who answers ``Eh, maybe because of heat conduction?'', is guessing the teacher's password.

An equation connects to experience only when you use it to predict the next measurement; even a picture of heat flowing counts if it leads to a test. The author proposes drilling students that only ``anticipation-controllers'' count, and lists the questions a student could ask to turn ``heat conduction?'' into experiments. Otherwise students get stuck on passwords, as the author did at nine and as humanity did, the post says, for thousands of years.

\responsehead

The post's central observation is sound. Students can give the approved answer without knowing what it means, and a system that grades answers encourages this. Its list of questions, which turns ``heat conduction?'' into a sequence of tests, is its most useful part.

Feynman made the same observation earlier, and the post does not say so. In ``O Americano, Outra Vez!'' Feynman describes Brazilian students who could define Brewster's angle and did not recognize it in the light off the bay: they ``had memorized everything, but they didn't know what anything meant.'' The post names two of Feynman's books, and uses his word ``wakalixes'' without saying where it comes from.

The post then claims more than it needs. It calls it a bad habit to think that the student who answers ``Waves!'' has made a true statement. That depends on a view of meaning which the post does not defend. On the view associated with Hilary Putnam and Tyler Burge, a speaker who uses a word as the experts do, and would accept their correction, means what the experts mean. The student's answer then says what the physicist's says, and the student's fault is a lack of understanding. The post's lesson, that understanding is what counts, holds on either view; the claim about truth is an extra step with no argument.

One passage contradicts the next. The post says that a student who ``just'' pictures heat flowing, and tests the picture with a match, is connected to experience. The next paragraph says that anyone not using the diffusion equation has that connection ``severed as though by a knife.''

Finally, the generalization is based on one case. The author at nine is the evidence that this ``is what happens by default,'' and that humanity ``stayed stuck in holes like this for thousands of years.'' The post gives no historical example.

In short: a sound point about memorized answers, made earlier by Feynman and uncredited, together with an unargued claim that the student's answer has no truth value, and extended to all of humanity from one child.
\end{afterword}
